\chapter{Introduction}\label{ch:intro}
In May 2017, the \textbf{WannaCry} ransomware outbreak dramatically underscored the fragility of modern operating systems. By exploiting an unpatched vulnerability in the Windows SMB protocol, WannaCry rapidly infected more than \textbf{230,000 computers across 150 countries}, crippling hospitals, businesses, and government services \citep{nist2017wannacry}. This was not an isolated incident; it highlighted a broader trend of escalating security failures in widely used operating systems. Linux, Windows, and macOS --- the foundational platforms for servers, desktops, and mobile devices --- have each been plagued by hundreds even thousands of new vulnerabilities every year. In 2024 alone, Microsoft Windows 10/11 (and various other server editions) had 4939 reported CVE-listed vulnerabilities, the mainline Linux kernel had 3889, and Apple’s macOS had 519 \citep{stackwatch2024stats}. Many of these flaws are critical, enabling remote code execution or privilege escalation by attackers. The steady drumbeat of patches and security updates for all major OSes attests to an uncomfortable reality: \textbf{modern operating systems remain fundamentally insecure despite their importance}.

While this is tolerable in general-purpose computing, it is unacceptable in safety- and security-critical systems --- such as those used in aircraft, medical devices, autonomous vehicles, and critical infrastructures --- where a single failure can lead to catastrophic outcomes \citep{butler1993-infeasibility, rushby1997-formal-methods}. These systems demand far stronger guarantees of correctness, isolation, and integrity than traditional OS architectures were ever designed to provide. In such domains, the operating system must not merely be \emph{``secure enough''}; it must be provably correct.

This has motivated decades of research into the building of a \textbf{Provably Secure Operating System (PSOS)} \citep{neumann1975psos, feiertag1979psos, neumann1980psos, neumann2003psos} --- systems for which critical properties can be mathematically verified. However, attempts to build such systems have historically been hindered by either impracticality or by excessive cost and complexity in the engineering and verification effort itself. That changed with the seL4 Microkernel \citep{klein2009seL4}, the first operating system kernel to be fully mathematically verified. seL4 provides strong spatial and temporal isolation, fine-grained access control via capabilities, and proofs of integrity, and confidentiality. The seL4 project demonstrates that formal verification can be applied to real-world systems of significant complexity.

The security guarantees provided by seL4 make it an attractive foundation for building a Provably Secure Operating System. And indeed, there have been some real-world deployments --- including autonomous aircraft \citep{Cofer_GBWPFPKKAHS_18}, and, more recently, autonomous vehicles \citep{Qu_24:sel4s}. However, more than a decade after seL4’s verification, we still haven’t seen a complete operating system stack built on top of it that is itself formally verified --- again owing to challenges in both the \textbf{engineering} and \textbf{verification} of the system.

To address these challenges, the seL4 Microkit \citep{microkit:url} and LionsOS \citep{heiser2025lionsos} were developed, demonstrating that it is possible to build a performant and practical system on top of seL4 without incurring significant engineering overhead. However, while these tools and design philosophies solve many of the engineering difficulties, the problem of verification beyond the kernel still requires further studies.

LionsOS aims to enable end-to-end verification of security and safety enforcement. Achieving this requires that components within its Trusted Computing Base (TCB) are not only functionally correct but formally verified down to the executable code. In seL4's verification story, the C compiler was removed from the TCB through translation validation \citep{sewell2013TV} --- a fragile and laborious process that imposes limitations on code complexity. A more robust approach is to use programming languages with well-defined semantics and a verified compiler that produce \emph{correct} binaries by construction.

This thesis aims to investigates three such programming languages and their compiler stacks for building a \textbf{provably secure} LionsOS without compromises to \textbf{performance} or \textbf{practicality}:
\begin{enumerate}
    \item \textbf{CompCert \citep{leroy2009compcert}:} a verified optimising compiler for the CLight dialect of C
    \item \textbf{CakeML \citep{kumar2014cakeml}:} a functional language with a end-to-end verified compiler
    \item \textbf{Pancake \citep{Pohjola2023pancake}:} a low-level systems language that builds upon CakeML's compiler backend.
\end{enumerate}

We will be providing a detailed, quantitative and qualitative analysis of each language system --- examining their suitability for different LionsOS components, the productivity impact on both developers and verifiers, and the performance of the generated code.

\section*{Structure Overview}
The remainder of this thesis will be structured as follows:
\todo[inline]{Add structure overview once the other sections are done}